\subsection{fhq-treap}
1. 值域分裂
\begin{lstlisting}
struct node{
    ll ls,rs;
    ll val,rnd;
    ll sz,num;
}t[1000010];ll cur=0;
ll newnode(ll val){
    t[++cur]={0,0,val,rand(),1,1};
    return cur;
}void upd(ll p){
    t[p].sz=t[t[p].ls].sz+t[t[p].rs].sz+t[p].num;
}void split(ll p,ll k,ll &x,ll &y){
    if(!p){x=y=0;return;}
    if(t[p].val<=k)x=p,split(t[p].rs,k,t[x].rs,y);
    else y=p,split(t[p].ls,k,x,t[y].ls);
    upd(p);
}ll merge(ll x,ll y){
    if(!x||!y)return x^y;
    if(t[x].rnd<t[y].rnd){
        t[x].rs=merge(t[x].rs,y);
        upd(x);return x;
    }else{
        t[y].ls=merge(x,t[y].ls);
        upd(y);return y;
    }
}
ll rt;
void init(){rt=cur=0;t[0]={};}
void insert(ll val){
    ll x,y,z;
    split(rt,val-1,x,y);split(y,val,y,z);
    if(y==0)y=newnode(val);else t[y].num++,t[y].sz++;
    rt=merge(merge(x,y),z);
}void del(ll val){
    ll x,y,z;
    split(rt,val-1,x,y);split(y,val,y,z);
    if(y==0){rt=merge(x,z);return;}t[y].num--,t[y].sz--;
    if(t[y].num==0)rt=merge(x,z);
    else rt=merge(merge(x,y),z);
}
ll rk(ll val){// 1 + (# of < val)
    ll x,y;split(rt,val-1,x,y);
    ll res=t[x].sz+1;rt=merge(x,y);
    return res;
}
ll findkth(ll p,ll k){
    while(1){
        if(t[t[p].ls].sz>=k)p=t[p].ls;
        else if(t[t[p].ls].sz+t[p].num>=k)return t[p].val;
        else k-=t[t[p].ls].sz+t[p].num,p=t[p].rs;
    }
}
ll pred(ll val){ // max_{x<val} x
    ll x,y;split(rt,val-1,x,y);
    ll res=findkth(x,t[x].sz);
    rt=merge(x,y);
    return res;
}ll succ(ll val){ // min_{x>val} x
    ll x,y;split(rt,val,x,y);
    ll res=findkth(y,1);
    rt=merge(x,y);
    return res;
}
\end{lstlisting}
2. 排名分裂
\begin{lstlisting}
ll n,m,cnt;
struct node{
    int ls,rs;
    ll val,rnd;
    int sz,num;
    int tg;// 打了 tag，可以不用。这里是区间 reverse 板子。
}t[100010];
#define ls(p) (t[(p)].ls)
#define rs(p) (t[(p)].rs)
int newnode(ll x){
    t[++cnt]={0,0,x,rand(),1,1,0};
    return cnt;
}
void upd(int p){t[p].sz=t[ls(p)].sz+t[rs(p)].sz+t[p].num;}
void pd(int p){
    if(t[p].tg==0)return;
    swap(ls(p),rs(p));
    if(ls(p))t[ls(p)].tg^=1;
    if(rs(p))t[rs(p)].tg^=1;
    t[p].tg=0;
}
void split(int p,int sz,int &x,int &y){
    if(!p){x=y=0;return;}pd(p);
    if(t[ls(p)].sz<sz)x=p,split(rs(p),sz-t[ls(p)].sz-t[p].num,rs(p),y);
    else y=p,split(ls(p),sz,x,ls(p));
    upd(p);
}int merge(int x,int y){
    if(!x||!y)return x+y;pd(x),pd(y);
    if(t[x].rnd<t[y].rnd){rs(x)=merge(rs(x),y);upd(x);return x;}
    else {ls(y)=merge(x,ls(y));upd(y);return y;}
}
void out(int p){
    if(!p)return;pd(p);
    out(ls(p)),cout<<t[p].val<<" ",out(rs(p));
}
ll rt;
void init(){rt=cnt=0;t[0]={};}
\end{lstlisting}

\paragraph{多测与注意} 两份代码独立选用，每组 init()。保留原 split、merge 和节点池；删除不回收编号。查询默认合法，val-1 要避开类型最小值；排名分裂版 num=1。
